\section{Theoretical Foundations}
\label{sec:theoretical}

% Full proofs and extended formal analysis relocated to Appendix~E.

We connect TSCG's operators to three foundational properties of causal autoregressive transformers.

\subsection{Causal Attention and Information Flow}

In a standard autoregressive transformer~\cite{vaswani2017attention}, attention at layer $\ell$ satisfies $\attn^{(\ell)}(i, j) = 0$ for $j > i$ (causal mask), creating asymmetric information flow: early tokens cannot access later ones.

\textbf{Implication for CFO.} If step $o_2$ depends on $o_1$ but $o_1$ appears \emph{after} $o_2$, the model must rely on parametric knowledge rather than direct attention.
CFO ensures $\pos(o_1) < \pos(o_2)$, guaranteeing prerequisites are causally accessible.

\textbf{Implication for CAS.} Total attention to position $i$ follows a U-shaped distribution~\cite{xiao2023efficient}: positions near 0 and $n$ receive disproportionate attention, while middle positions form an ``attention valley.''
CAS places high-fragility tools at positions~0 and~$n$.

\subsection{Attention Sink Exploitation}

The attention sink phenomenon~\cite{xiao2023efficient}---position~0 receives $\attn(i, 0) > 1/i$ for most $i$ and layers---provides a natural amplification mechanism.

\textbf{Implication for CFL.} Placing the output constraint at position~0 ensures it receives elevated attention from every subsequent token.
Combined with CCP/SAD-F at position~$n$ (recency bias), this creates a ``bookend'' strategy: critical information occupies both extremes of the attention distribution.

\subsection{BPE Non-Monotonicity}

\begin{theorem}[Tokenization Non-Monotonicity]
\label{thm:non-monotonicity}
For BPE tokenizer $\mathcal{T}$ and strings $s_1, s_2$ with $|s_1|_\text{chars} < |s_2|_\text{chars}$, it is not necessarily the case that $|\tok(s_1)| < |\tok(s_2)|$.
\end{theorem}

TAS exploits this by selecting surface forms aligned with learned BPE merges, achieving token reduction without semantic change.
Proof in Appendix~E.

\subsection{Attention Dilution and SDM}

\begin{proposition}[SDM Improves Effective Attention]
\label{prop:sdm-attention}
Removing $k$ filler tokens from a prompt of length $n$ increases average effective attention per semantic atom by at least $n/(n-k)$.
\end{proposition}

This follows from softmax normalization: filler tokens compete for attention weight with semantic tokens.
SDM removes this competition, providing formal justification beyond token cost savings.

\subsection{Fragility and Causal Accessibility}

\begin{definition}[Causal Accessibility]
\label{def:causal-accessibility}
$\mathcal{A}(a) = \frac{1}{L} \sum_{\ell=1}^{L} \attn^{(\ell)}(n, i)$: average attention from the generation position to atom $a$ at position $i$.
\end{definition}

\begin{definition}[Fragility]
\label{def:fragility-gap}
An atom $a$ is fragile when importance exceeds accessibility: $\fragility(a) \propto \text{importance}(a) - \mathcal{A}(a)$.
\end{definition}

High-importance, low-accessibility atoms are the targets for CAS reordering and SAD-F duplication.
Budget-constrained duplication is preferable to naive repetition because duplicating all atoms would increase length and dilute attention (Proposition~\ref{prop:sdm-attention}).
